Large language model (LLM) agents are increasingly deployed on data-science tasks, yet existing benchmarks report only aggregate success rates, masking \emph{where} and \emph{why} agents fail. We introduce \textbf{AgentFail}, a diagnostic framework that decomposes every agent failure into five stages---Analytical-Plan, Code-Generation, Runtime, Output-Mismatch, and Answer-Error---and distinguishes \emph{silent failures} (code executes without error but yields a wrong answer) from \emph{loud failures} (code crashes). We propose \emph{oracle-guided repair} to test repairability, with no-op negative controls. We evaluate five frontier LLMs across 5{,}984 runs on three task suites: 95 synthetic trap tasks, 23 real UCI datasets, and 200 real DSBench tasks. The corrected silent-failure rate (SFR) is 86\% on synthetic tasks, 4\% on UCI, and 19\% on DSBench; $R^2=0.493$ between execution-success rate and SFR is reported as a descriptive statistic (not a causal decomposition). Recovery rate is 0\% across all models. A failure-aware verifier helps two of five models (McNemar $p<0.001$). Oracle repair matches the no-op baseline (64\%), establishing an upper bound on reparability rather than step-level causal attribution. Cohen's $\kappa=0.435$ (moderate, LLM-as-judge). All tables are auto-generated from a single manifest with assertion checks. We release all code, data, and 5{,}984 annotated runs.
